% ============================================================
% Preámbulo canónico para una ficha de auditoría
%
% Cópialo tal cual y compila con pdfLaTeX, DOS pasadas.
% Rellena \hypersetup con tus propios datos antes de entregar.
%
% Cada bloque comentado documenta una trampa que ya costó una
% iteración; no las deshagas sin motivo. Ajusta sólo la geometría y
% el interlineado si la consigna de la actividad pide otra cosa.
% ============================================================
\documentclass[12pt,letterpaper]{article}

% --- Codificación e idioma ---
% es-nodecimaldot: sin esta opción, babel-spanish imprime los decimales
% con coma en modo matemático.
\usepackage[T1]{fontenc}
\usepackage[utf8]{inputenc}
\usepackage[spanish,es-nodecimaldot]{babel}

% --- Tipografía Times ---
% Times New Roman es propietaria de Microsoft y Overleaf no la distribuye.
% newtx es el clon con las mismas métricas. NO lo cambies por mathptmx:
% ése reclama la fuente script rsfs10 y falla en instalaciones parciales.
\usepackage{newtxtext,newtxmath}
\usepackage{microtype}
\setlength{\emergencystretch}{3em}

% --- Geometría ---
\usepackage[letterpaper,left=1.1in,right=1.1in,top=0.95in,bottom=0.95in]{geometry}

% --- Interlineado y párrafo ---
% 1.12 mantiene una ficha densa legible. Si la consigna exige 1.5,
% usa \onehalfspacing y recalcula el presupuesto de páginas.
\usepackage{setspace}
\setstretch{1.12}
\setlength{\parindent}{0pt}
\setlength{\parskip}{4pt plus 2pt minus 1pt}

% --- Folio centrado al pie ---
\usepackage{fancyhdr}
\pagestyle{fancy}
\fancyhf{}
\fancyfoot[C]{\thepage}
\renewcommand{\headrulewidth}{0pt}
\fancypagestyle{plain}{\fancyhf{}\fancyfoot[C]{\thepage}\renewcommand{\headrulewidth}{0pt}}

% --- Paleta de marca (cámbiala por la tuya) ---
% Los mismos valores que vp_model/palette.py. No inventes otros azules.
\usepackage[table]{xcolor}
\definecolor{marcablue}{HTML}{003CA6}
\definecolor{marcayellow}{HTML}{FFD600}
\definecolor{marcagray}{HTML}{555559}
\definecolor{marcastripe}{HTML}{F4F6FB}

% --- Gráficos ---
\usepackage{graphicx}
\usepackage{float}                % [H] para figuras pegadas a su párrafo
\usepackage{tikz}                 % recorte circular del retrato de portada
\usepackage{caption}
\captionsetup{font={footnotesize,color=marcagray},labelfont={bf,color=marcagray},
              justification=justified,format=plain,skip=5pt}

% --- Tablas ---
% pifont da la palomita de las listas de verificación. Con babel-spanish,
% el \checkmark de amssymb reclama la fuente script rsfs10, que no toda
% instalación trae; \ding{51} es un glifo de texto y no la pide.
\usepackage{array}
\usepackage{tabularx}
\usepackage{booktabs}
\usepackage{enumitem}
\usepackage{amsmath}
\usepackage{pifont}
\renewcommand{\arraystretch}{1.06}
\newcolumntype{L}{>{\raggedright\arraybackslash}X}

% Ancho de la columna de etiquetas en los bloques de clave y valor.
\newlength{\etiq}\setlength{\etiq}{0.235\textwidth}
\newcolumntype{E}{>{\bfseries\raggedright\arraybackslash}p{\etiq}}

% Las tablas de ficha son cajas indivisibles: al interlineado del cuerpo
% una sola no cabe en el hueco restante y arrastra media página en blanco.
% De ahí \canvasfont, que las compacta.
%
% ⚠️ NO envuelvas tabularx en un \newenvironment: explora su cuerpo buscando
% literalmente \end{tabularx} y revienta con «Argument of \TX@get@body has
% an extra }». Cada tabla se escribe completa:
%
%   \begingroup\canvasfont
%   \begin{tabularx}{\textwidth}{@{}E L@{}}
%   \toprule
%   Etiqueta & Valor \\
%   \bottomrule
%   \end{tabularx}
%   \endgroup
\newcommand{\canvasfont}{\small\setstretch{1.02}}

% Nota al pie de un bloque: comentario del autor, no parte de la plantilla.
\newcommand{\nota}[1]{\par\addvspace{2pt}{\footnotesize\itshape\textcolor{marcagray}{#1}}\par}

% --- Secciones sin numerar (el número va embebido en el rótulo) ---
\setcounter{secnumdepth}{-1}
\usepackage{titlesec}
\titleformat{\section}
  {\normalfont\fontsize{12.5}{15}\bfseries\color{marcablue}}{}{0pt}{}
\titlespacing*{\section}{0pt}{7pt plus 2pt minus 2pt}{2pt}
\titleformat{\subsection}
  {\normalfont\fontsize{12}{14}\bfseries}{}{0pt}{}
\titlespacing*{\subsection}{0pt}{8pt plus 2pt minus 2pt}{1pt}

% Bloque numerado de la plantilla: filete azul, número y título en una línea.
% \texorpdfstring evita que el \quad llegue al marcador del PDF y dispare
% el aviso «Token not allowed in a PDF string» de hyperref.
\newcommand{\bloque}[2]{%
  \par\addvspace{7pt plus 3pt minus 2pt}%
  \noindent\textcolor{marcablue}{\rule{\textwidth}{0.9pt}}%
  \par\vspace{1pt}%
  \section{\texorpdfstring{#1\quad #2}{#1 #2}}%
  \nopagebreak%
}

% --- Bibliografía sin su propio encabezado: el rótulo lo pone \bloque ---
% Sin esto, thebibliography imprime su \section* y queda un título duplicado.
\makeatletter
\renewenvironment{thebibliography}[1]
  {\list{\@biblabel{\@arabic\c@enumiv}}%
        {\settowidth\labelwidth{\@biblabel{#1}}%
         \leftmargin\labelwidth \advance\leftmargin\labelsep
         \usecounter{enumiv}\let\p@enumiv\@empty
         \renewcommand\theenumiv{\@arabic\c@enumiv}}%
   \sloppy\clubpenalty4000 \@clubpenalty\clubpenalty \widowpenalty4000%
   \sfcode`\.\@m}
  {\def\@noitemerr{\@latex@warning{Empty `thebibliography' environment}}%
   \endlist}
\makeatother

% --- Enlaces ---
\usepackage{xurl}
\usepackage[hidelinks,breaklinks=true]{hyperref}
\hypersetup{
  pdftitle={},
  pdfauthor={},
  pdfsubject={},
  pdflang={es-MX}
}

% ============================================================
% RECORDATORIOS DE CUERPO
%
% · \begin{thebibliography}{9} basta hasta 9 referencias; con 10 o más
%   usa {99} o la etiqueta [10] desborda la sangría francesa.
% · El orden de los \bibitem ES la numeración IEEE. Ordénalos por primera
%   aparición en el texto y verifícalo por script, nunca a ojo.
% · \% siempre en modo texto: 58\,\%, jamás dentro de $…$.
% · Nunca «shorten <=» ni «shorten >=» en TikZ: babel-spanish activa < y >
%   como caracteres especiales y rompe el parser de claves.
% · Toda figura necesita \label y al menos un \ref, y la mención va antes.
% ============================================================
